\documentclass[11pt]{article}
\usepackage[utf8]{inputenc}
\usepackage[margin=1in]{geometry}
\usepackage{booktabs}
\usepackage{hyperref}
\usepackage{enumitem}
\usepackage{amsmath}

\title{Alumni Community Platform: Feature Development Plan}
\author{Patricio Jose Najeal, AI-Generated (ChatGPT, Gemini)}
\date{March 17, 2026}

\begin{document}

\maketitle

\section{System Overview}
The platform is a moderated alumni community that integrates microblogging, event management, and a friction-less donation system. The architecture balances casual user engagement with high-reliability administrative oversight.

\section{Module 1: Microblogging \& Social Feed}
\begin{itemize}
    \item \textbf{Post Creation:} All registered users are authorized to create posts.
    \item \textbf{Infinite Scroll:} Recommended for the public feed to mimic modern social media behavior and encourage casual browsing.
    \item \textbf{Batch Loading:} Data is fetched in batches (e.g., 20--30 posts) to prevent excessive database queries.
    \item \textbf{In-App Notifications:} Used for social interactions (replies, mentions, approvals) to keep discussions inside the platform and avoid email fatigue.
\end{itemize}

\section{Module 2: Event Management \& RSVP}
\begin{itemize}
    \item \textbf{Verification:} Organizers must be privileged users who have completed an identity verification process.
    \item \textbf{Email Communications:} Logistics, reminders, and RSVP confirmations are sent via email to ensure visibility and calendar integration.
    \item \textbf{RSVP Mechanics:} Supports "Going," and "Not Going" statuses with strict configurable deadlines.
    \item \textbf{Attendance Management:} Automated waitlisting triggers when capacity limits are reached; the system notifies the next user in line if a slot opens.
\end{itemize}

\pagebreak

\section{Module 3: Moderation \& NLP Layer}
\begin{itemize}
    \item \textbf{NLP Screening:} An automated layer flags harassment, spam, or offensive language for review.
    \item \textbf{Admin Pipeline:} All flagged content and new events must be approved by an administrator before public visibility.
    \item \textbf{Paginated Workflow:} Administrative tasks (review queues, user verification) use pagination for stable, deterministic ordering.
\end{itemize}

\section{Module 4: Donation System}
\begin{itemize}
    \item \textbf{Low-Friction Access:} Guest donations are permitted without account creation to maximize contributions.
    \item \textbf{Profile Linking:} Logged-in users have donations automatically linked; guest donors can "claim" donations later by verifying their email.
    \item \textbf{Anonymity:} Supports public anonymity while maintaining internal metadata (Transaction IDs, IP logs) for accounting and fraud prevention.
\end{itemize}

\newpage
\section{Database Schema: RSVP \& Waitlist System}
The following relational schema is designed to handle event capacity, chronological waitlisting, and attendance tracking.

\subsection{Events Table}
Defines the constraints for attendance.
\begin{itemize}
    \item \texttt{event\_id} (UUID, PK)
    \item \texttt{organizer\_id} (UUID, FK)
    \item \texttt{max\_capacity} (INT) --- Maximum number of allowed "Going" statuses.
    \item \texttt{rsvp\_deadline} (DATETIME) --- Cutoff for user responses.
    \item \texttt{is\_approved} (BOOLEAN) --- Administrative toggle.
\end{itemize}

\subsection{RSVPs Table}
Tracks active participation responses.
\begin{itemize}
    \item \texttt{rsvp\_id} (UUID, PK)
    \item \texttt{event\_id} (UUID, FK)
    \item \texttt{user\_id} (UUID, FK)
    \item \texttt{status} (ENUM: 'Going', 'Not Going')
    \item \texttt{guest\_count} (INT) --- Default is 0.
    \item \texttt{updated\_at} (TIMESTAMP)
\end{itemize}

\subsection{Waitlist Table}
Manages the queue for over-capacity events.
\begin{itemize}
    \item \texttt{waitlist\_id} (UUID, PK)
    \item \texttt{event\_id} (UUID, FK)
    \item \texttt{user\_id} (UUID, FK)
    \item \texttt{queue\_position} (SERIAL) --- Used for FIFO (First-In, First-Out) logic.
    \item \texttt{created\_at} (TIMESTAMP)
\end{itemize}

\subsection{Capacity Logic}
To determine if a user should be moved from the waitlist to active status, the system evaluates:
$$ \text{Total Confirmed} = \sum (\text{RSVPs where status} = \text{'Going'} + \text{guest\_count}) $$
If $\text{Total Confirmed} < \text{max\_capacity}$, the user at $\min(\text{queue\_position})$ is promoted.

\end{document}